
\documentclass[11pt]{article}
\usepackage[T1]{fontenc}
\usepackage[utf8]{inputenc}
\usepackage[margin=1in]{geometry}
\usepackage{amsmath,amssymb,amsthm,mathtools}
\usepackage[hidelinks]{hyperref}
\usepackage{enumitem}
\usepackage{array}
\usepackage{longtable}

\theoremstyle{definition}
\newtheorem{definition}{Definition}[section]
\newtheorem{convention}[definition]{Convention}
\newtheorem{implementation}[definition]{Implementation choice}
\newtheorem{formalnote}[definition]{Formalization note}
\newtheorem{checklist}[definition]{Checklist}
\newtheorem{contract}[definition]{Lean contract}

\theoremstyle{plain}
\newtheorem{theorem}[definition]{Theorem}
\newtheorem{lemma}[definition]{Lemma}
\newtheorem{proposition}[definition]{Proposition}
\newtheorem{corollary}[definition]{Corollary}

\theoremstyle{remark}
\newtheorem{remark}[definition]{Remark}

\newcommand{\N}{\mathbb{N}}
\newcommand{\Z}{\mathbb{Z}}
\newcommand{\R}{\mathbb{R}}
\newcommand{\Finset}{\operatorname{Finset}}
\newcommand{\Set}{\operatorname{Set}}
\newcommand{\Sgn}{\operatorname{Sign3}}
\newcommand{\sgn}{\operatorname{sgn}}
\newcommand{\opp}{\operatorname{opp}}
\newcommand{\defect}{\operatorname{defect}}
\newcommand{\nabs}[1]{\left|#1\right|_{\N}}
\newcommand{\card}[1]{\left|#1\right|}
\newcommand{\range}{\operatorname{range}}
\newcommand{\SumN}{\operatorname{sumN}}
\newcommand{\CoeffOn}{\operatorname{CoeffOn}}
\newcommand{\Grade}{\operatorname{grade}}
\newcommand{\Val}{\operatorname{value}}
\newcommand{\Diff}{\operatorname{Diff}}
\newcommand{\AdmF}{\operatorname{AdmFin}}
\newcommand{\AdmS}{\operatorname{AdmSet}}
\newcommand{\PaperAdmF}{\operatorname{PaperAdmFin}}
\newcommand{\PaperAdmS}{\operatorname{PaperAdmSet}}
\newcommand{\celt}{\operatorname{c}}
\newcommand{\Cblock}{\operatorname{C}}
\newcommand{\Qfun}{\operatorname{Q}}
\newcommand{\CDiff}{\operatorname{CDiff}}
\newcommand{\LGrade}{\operatorname{localGrade}}
\newcommand{\LIndex}{\operatorname{localIndex}}
\newcommand{\LVal}{\operatorname{localValue}}
\newcommand{\DTotal}{\operatorname{defectTotal}}
\newcommand{\WDTotal}{\operatorname{weightedDefectTotal}}
\newcommand{\USmall}{\operatorname{USmall}}
\newcommand{\WSmall}{\operatorname{WSmall}}
\newcommand{\HSmall}{\operatorname{HSmall}}
\newcommand{\OldSmall}{\operatorname{OldSmall}}
\newcommand{\NewSmall}{\operatorname{NewSmall}}
\newcommand{\CarryState}{\operatorname{CarryState}}
\newcommand{\LegalStep}{\operatorname{LegalStep}}
\newcommand{\StepOut}{\operatorname{StepOut}}
\newcommand{\block}{\operatorname{block}}
\newcommand{\elem}{\operatorname{elem}}
\newcommand{\owner}{\operatorname{owner}}
\newcommand{\offset}{\operatorname{offset}}
\newcommand{\gap}{\operatorname{gap}}
\newcommand{\seq}{\operatorname{seq}}
\newcommand{\thetaC}{\theta}
\newcommand{\tauC}{\tau}
\newcommand{\fexp}{f}
\newcommand{\rpow}{\operatorname{rpow}}
\newcommand{\atTop}{\operatorname{atTop}}
\newcommand{\nhds}{\mathcal{N}}
\newcommand{\ceil}{\operatorname{ceil}}
\newcommand{\eps}{\varepsilon}
\newcommand{\Geom}{\operatorname{Geom}}
\newcommand{\Pfun}{\operatorname{P}}
\newcommand{\Scale}{\operatorname{scale}}
\newcommand{\ZLtQuarter}{\operatorname{ZLtQuarter}}
\newcommand{\ZLtHalf}{\operatorname{ZLtHalf}}
\newcommand{\Shift}{\operatorname{shift}}
\newcommand{\AdmSeq}{\operatorname{AdmSeq}}
\newcommand{\PrefixRange}{\operatorname{PrefixRange}}
\newcommand{\GapAt}{\operatorname{GapAt}}
\newcommand{\AnalysisReady}{\operatorname{AnalysisReady}}

\title{Lean-First Blueprint for the Admissible Carry Construction\\
Version 10: two-sided local core, sequence-first admissibility, max-one scale}
\author{Prepared for formalization of the May 2026 note}
\date{May 2026}

\begin{document}
\maketitle

\begin{abstract}
This document is a further Lean-oriented rewrite of the admissible carry construction in the attached May 2026 note.  Version 10 keeps the previous single proof route, range-sum local blocks, cached carry states, and block-owner enumeration, but removes three more implementation pain points.  First, the preferred local carry theorem is stated with two-sided integer inequalities \(\ZLtQuarter\) and \(\ZLtHalf\), so the hardest arithmetic proof no longer sees \(\operatorname{Int.natAbs}\).  Second, admissibility is proved first for the constructed sequence via finite prefixes, and only then converted to the paper's set formulation.  Third, the asymptotic layer uses the total scale \(\Scale(j)=\max\{1,N_j\}\) and \(\Pfun(j)=M_j/\Scale(j)^\tau\), avoiding shifted denominators and the recurring \(j\ge1\) split.  The mathematical construction remains aligned with the attached note; only the formalization interface is changed.
\end{abstract}

\tableofcontents

\section{Difficulty rating}

\begin{formalnote}[Current rating]
From the Version 9 blueprint I would rate the Lean formalization difficulty at about
\[
        \boxed{4.5/10}
\]
for an experienced Lean user.  Version 10 should lower the practical difficulty to roughly
\[
        \boxed{4.25/10}.
\]
The finite combinatorial core is now about \(3.5\)--\(4/10\), because the local-carry arithmetic is shielded from \(\operatorname{Int.natAbs}\) and the closure induction uses cached records.  The full theorem, including exact enumeration, concrete ceiling estimates, real powers, and filters, is about \(4.25\)--\(4.5/10\).  For someone not already comfortable with \verb|Finset| sums, integer casts, \verb|omega|/\verb|nlinarith|, \verb|Nat.find|, records, and \verb|Real.rpow|, the same blueprint is closer to \(5.5\)--\(6/10\).
\end{formalnote}

\begin{formalnote}[Why the rating improves]
The main improvement is not a new mathematical argument; it is the removal of formalization forks.  Earlier versions contained several ways to prove local carry: subset pairs, \(\operatorname{Fin} m\) arrays, signed coefficients, and paper-style auxiliary integers.  Version 10 keeps exactly one route and adds two boundary wrappers: a sequence-level admissibility wrapper and a max-one asymptotic scale.  The finite core still follows this route:
\[
\text{range coefficients}\Rightarrow
\text{defect totals}\Rightarrow
\text{positive shift contradiction}\Rightarrow
\text{local carry}\Rightarrow
\text{closure}.
\]
This makes the theorem dependency graph substantially smaller and reduces coercion drift.
\end{formalnote}

\begin{formalnote}[Remaining hard spots]
The remaining nontrivial work is concentrated in five places.
\begin{enumerate}[label=(\arabic*), leftmargin=*]
\item The positive shifted local-carry contradiction, because it mixes \(\N\)-valued defect sums with \(\Z\)-valued index sums.
\item The block-dilation bridge \(\Diff(MC_k,t,x)\Rightarrow x=M y\) and \(\CDiff(k,t,y)\), because it involves images of finite sets.
\item The closure proof, because it decomposes one signed difference over a disjoint union and then recombines two carries.
\item The exact increasing enumeration, because it requires a block-owner function and one boundary case at the end of each block.
\item The concrete asymptotic wrapper, because it uses ceilings, \(\operatorname{Real.rpow}\), and filter limits.
\end{enumerate}
All five are now isolated behind small theorem statements.
\end{formalnote}

\section{Alignment with the attached note}

\begin{convention}[Zero-indexed stage notation]
The attached note uses paper stages
\[
        A_0=\varnothing,
        \qquad M_1=1,
        \qquad N_0=0,
\]
and for \(p\ge1\)
\[
        k_p=\left\lceil 2(1+N_{p-1})^{1+\sqrt2}\right\rceil,
        \quad B_p=M_pC_{k_p},
        \quad A_p=A_{p-1}\cup B_p,
        \quad M_{p+1}=M_pQ_{k_p}.
\]
In Lean, use zero-indexed stages
\[
        F_0=\varnothing,
        \qquad M_0=1,
        \qquad N_0=0,
\]
and for \(j:\N\)
\[
        k_j=\left\lceil 2(1+N_j)^{1+\sqrt2}\right\rceil,
        \quad B_j=M_jC_{k_j},
        \quad F_{j+1}=F_j\cup B_j,
        \quad M_{j+1}=M_jQ_{k_j},
        \quad N_{j+1}=N_j+k_j.
\]
The dictionary is
\[
\begin{array}{c|c}
\text{paper note} & \text{Lean blueprint}\\ \hline
A_p & F_p\\
N_p & N_p\\
B_p & B_{p-1}\quad(p\ge1)\\
k_p & k_{p-1}\quad(p\ge1)\\
M_p & M_{p-1}\quad(p\ge1).
\end{array}
\]
Thus the paper's first block \(B_1=\{1,4\}\) is Lean's block \(B_0\).
\end{convention}

\begin{formalnote}[What Version 10 keeps from Version 9 and adds]
Version 10 makes the following formalization decisions binding.
\begin{enumerate}[label=(\arabic*), leftmargin=*]
\item Local block sums are over \(\operatorname{Finset.range} m\), not over \(\operatorname{Fin} m\).  This exposes the bound \(i<m\) directly and avoids most \(\operatorname{Fin.val}\) casts.
\item General finite differences use a \(\CoeffOn(A)\) record with a support-zero field.  Local block differences use the simpler object \(\epsilon:\N\to\operatorname{Sign3}\), because only values on \(\operatorname{range} m\) are ever summed.
\item The older paper-style local-carry route using a free integer \(a\) is no longer part of the main dependency chain.  The Nat-valued defect total \(R=\DTotal(m,\epsilon)\) is the canonical replacement.
\item Carry states cache \(\sigma=\SumN(A)\).  Smallness lemmas use \(\sigma\) and do not unfold \(\SumN(A)\) except at construction time.
\item The increasing enumeration is not defined by sorting the infinite set.  It is defined by a block owner and an offset, and then proved to enumerate the same set.
\end{enumerate}
\end{formalnote}


\section{Version 10 additional implementation contracts}\label{sec:v10-contracts}

\begin{formalnote}[Three additional binding choices]
Version 10 makes the following choices binding for the Lean implementation.
\begin{enumerate}[label=(\arabic*), leftmargin=*]
\item The local carry theorem used by closure is a wrapper around a preferred two-sided integer theorem.  Only the wrapper mentions \(\operatorname{Int.natAbs}\).
\item The infinite construction first proves a sequence-level admissibility predicate \(\AdmSeq(\seq)\).  The paper set statement for \(\Set.range(\seq)\) is a final wrapper.
\item The real-asymptotic proof uses \(\Scale(j)=\max\{1,N_j\}\), so every normalized denominator is positive and no lemma needs a recurring \(j\ge1\) hypothesis merely to avoid \(N_0=0\).
\end{enumerate}
\end{formalnote}

\begin{definition}[Two-sided local smallness]\label{def:v10-zsmall}
For \(m:\N\) and \(x:\Z\), define
\[
        \ZLtQuarter(m,x)\Longleftrightarrow -(m:\Z)<4x\ \wedge\ 4x<(m:\Z),
\]
\[
        \ZLtHalf(m,x)\Longleftrightarrow -(m:\Z)<2x\ \wedge\ 2x<(m:\Z).
\]
The public predicates \(\USmall\) and \(\WSmall\) remain useful for quotient extraction, but the local-carry proof should immediately convert them to \(\ZLtQuarter\) and \(\ZLtHalf\).
\end{definition}

\begin{lemma}[NatAbs firewall for local carry]\label{lem:v10-natabs-firewall}
For all \(m,u,w\), prove
\[
        \USmall(m,u)\to\ZLtQuarter(m,u),
        \qquad
        \WSmall(m,w)\to\ZLtHalf(m,w),
\]
\[
        \ZLtQuarter(m,u)\leftrightarrow \ZLtQuarter(m,-u),
        \qquad
        \ZLtHalf(m,w)\leftrightarrow \ZLtHalf(m,-w).
\]
These are the only local-carry lemmas that should unfold or reason directly about \(\operatorname{Int.natAbs}\).
\end{lemma}

\begin{definition}[Shift normal form]\label{def:v10-shift}
For \(m:\N\), \(\epsilon:\N\to\Sgn\), and \(w:\Z\), define
\[
        \Shift(m,\epsilon,w)=(m:\Z)w-\LIndex(m,\epsilon).
\]
The quotient identity should be stated in the exact form
\[
        \LVal(m,\epsilon)=\Qfun(m)w-u
        \Longrightarrow
        \LGrade(m,\epsilon)=(w-u)+((m:\Z)+1)\Shift(m,\epsilon,w).
\]
This is a \verb|ring| lemma, not a modular arithmetic lemma.
\end{definition}

\begin{theorem}[Preferred two-sided local carry theorem]\label{thm:v10-local-carry-zsmall}
Let \(2\le m\), let \(\epsilon:\N\to\Sgn\), and suppose
\[
        \LVal(m,\epsilon)=\Qfun(m)w-u,
        \qquad
        \ZLtQuarter(m,u),
        \qquad
        \ZLtHalf(m,w).
\]
Then
\[
        \LGrade(m,\epsilon)=w-u.
\]
The older Theorem~\ref{thm:v9-local-carry-range} should be derived as a wrapper using Lemma~\ref{lem:v10-natabs-firewall} and the construction fact \(2\le k_j\).  In practice, every local-carry call in closure has \(m=k_j\), so the stronger lower bound is available.
\end{theorem}

\begin{lemma}[Three-shift arithmetic core]\label{lem:v10-three-shift-core}
Let \(2\le m\), and let \(G,T,s:\Z\).  If
\[
        -(m:\Z)\le G\le (m:\Z),
        \qquad
        -3(m:\Z)<4T,
        \qquad
        4T<3(m:\Z),
\]
\[
        G=T+((m:\Z)+1)s,
\]
then
\[
        s=-1\quad\vee\quad s=0\quad\vee\quad s=1.
\]
This theorem contains no finite sums, no signs, and no block definitions.  It should be an \verb|omega| lemma.  The local-carry proof applies it with \(G=\LGrade(m,\epsilon)\), \(T=w-u\), and \(s=\Shift(m,\epsilon,w)\).
\end{lemma}

\begin{lemma}[Two-sided smallness bounds \(w-u\)]\label{lem:v10-t-bound}
If \(\ZLtQuarter(m,u)\) and \(\ZLtHalf(m,w)\), then
\[
        -3(m:\Z)<4(w-u),
        \qquad
        4(w-u)<3(m:\Z).
\]
After unfolding Definition~\ref{def:v10-zsmall}, this is pure integer arithmetic.
\end{lemma}

\begin{definition}[Sequence-level admissibility]\label{def:v10-admseq}
For a sequence \(a:\N\to\N\), define
\[
        \AdmSeq(a)\Longleftrightarrow
        \forall U,V:\Finset\,\N,
        \sum_{i\in U}a_i=\sum_{i\in V}a_i\Longrightarrow \card U=\card V.
\]
This is the preferred infinite admissibility target in Lean.  It quantifies over finite sets of indices rather than finite subsets of \(\Set.range(a)\).
\end{definition}

\begin{definition}[Prefix-range invariant]\label{def:v10-prefix-range}
For the zero-based sequence \(\seq\), define
\[
        \PrefixRange(j)\Longleftrightarrow
        F_j=(\Finset.range\,N_j).\operatorname{image}(\seq).
\]
The stage induction should prove \(\PrefixRange(j)\) along with carry-goodness.  This makes the proof of \(\AdmSeq(\seq)\) a finite-prefix argument.
\end{definition}

\begin{theorem}[Admissible prefixes imply sequence admissibility]\label{thm:v10-admseq-from-prefixes}
Assume \(\seq\) is strictly increasing, \(N_j\to\infty\), \(\PrefixRange(j)\) for every \(j\), and every finite stage \(F_j\) is admissible.  Then
\[
        \AdmSeq(\seq).
\]
Given finite index sets \(U,V\), choose \(J\) with \(U\cup V\subseteq\Finset.range(N_J)\).  The image sets lie in \(F_J\); strict monotonicity preserves cardinals under image; finite-stage admissibility gives the result.
\end{theorem}

\begin{theorem}[Paper set wrapper]\label{thm:v10-paper-set-wrapper}
If \(\seq\) is strictly increasing and \(\AdmSeq(\seq)\), then \(\Set.range(\seq)\) is admissible in the paper sense.  This should be the only theorem that converts finite subsets of a range back to finite sets of indices.
\end{theorem}

\begin{definition}[Additive gap target]\label{def:v10-gapat}
Define
\[
        \GapAt(n)=M_{\owner(n)}(k_{\owner(n)}+1).
\]
The primary exact gap theorem should be
\[
        \seq(n+1)=\seq(n)+\GapAt(n).
\]
Only after this additive equality is proved should one derive the subtraction form
\[
        \seq(n+1)-\seq(n)=\GapAt(n).
\]
This isolates all natural-subtraction goals in one wrapper.
\end{definition}

\begin{definition}[Total asymptotic scale]\label{def:v10-scale}
Define
\[
        \Scale(j)=\max\{1,N_j\},
        \qquad
        \Pfun(j)=\frac{(M_j:\R)}{\rpow((\Scale(j):\R),\tauC)}.
\]
The concrete ceiling file should prove
\[
        2\rpow((\Scale(j):\R),\thetaC)\le(k_j:\R),
        \qquad
        (k_j+1:\R)\le C_0\rpow((\Scale(j):\R),\thetaC)
\]
for every \(j\), not just for \(j\ge1\).  A slightly larger constant \(C_0\) absorbs the base case.
\end{definition}

\begin{definition}[Analysis-ready record]\label{def:v10-analysis-ready}
The final real-analysis file should prove one record \(\AnalysisReady\) containing:
\begin{enumerate}[label=(\arabic*), leftmargin=*]
\item the additive gap theorem for \(\seq\);
\item \(\owner(n)\to\atTop\);
\item \(\Pfun(j)\to0\);
\item an eventual bound
\[
        0\le\frac{(\GapAt(n):\R)}{\rpow((n+1:\R),\fexp)}
        \le C\Pfun(\owner(n)).
\]
\end{enumerate}
The theorem \(\AnalysisReady\to\) little-oh should be a one-time filter wrapper and should not unfold the construction.
\end{definition}

\section{File separation and import barriers}

\begin{contract}[Suggested file order]
The proof should be split as follows.
\begin{center}
\begin{tabular}{>{\raggedright\arraybackslash}p{0.28\linewidth} >{\raggedright\arraybackslash}p{0.62\linewidth}}
File & Contents \\ \hline
\verb|Signs| & \(\operatorname{Sign3}\), \(\sgn\), \(\opp\), \(\defect\), sign simp lemmas.\\
\verb|CoeffDiff| & \(\CoeffOn\), \(\Grade\), \(\Val\), \(\Diff\), bridges to subset pairs, finite-sum bounds.\\
\verb|LocalRange| & \(\celt\), \(\Cblock\), \(\Qfun\), range-sum local grade/index/value, \(\CDiff\).\\
\verb|LocalCarry| & defect totals, three-case reduction, positive and negative shift impossibility, local carry theorem.\\
\verb|Closure| & integer smallness predicates, cached carry states, one-step closure.\\
\verb|Stage| & zero-indexed recursive construction from abstract legal block sizes.\\
\verb|Enumeration| & block owner, offset, exact sequence formula, exact gap formula.\\
\verb|AsymptoticAbstract| & geometric-decay lemma and block-ratio-to-sequence-ratio lemma.\\
\verb|ConcreteCeil| & ceiling choice \(k_j=\lceil2(1+N_j)^\theta\rceil\), real-power estimates, final theorem.
\end{tabular}
\end{center}
The first seven files should not import real analysis or filters.  The only files involving \(\R\), \(\operatorname{Real.rpow}\), ceilings as real estimates, or \(\atTop\) should be \verb|AsymptoticAbstract| and \verb|ConcreteCeil|.
\end{contract}

\begin{formalnote}[Local simp sets]
Use local simp collections rather than global unfolding.
\begin{enumerate}[label=(\roman*), leftmargin=*]
\item \verb|sign_simps|: definitions and case lemmas for \(\sgn\), \(\opp\), and \(\defect\).
\item \verb|range_local_simps|: \(\celt\), \(\Qfun\), \(\LGrade\), \(\LIndex\), \(\LVal\), \(\DTotal\), and \(\WDTotal\), but only inside the local-carry file.
\item \verb|state_simps|: recursive equations for \(F\), \(M\), \(N\), \(k\), \(B\), and cached sums, used only in the stage and enumeration files.
\end{enumerate}
Do not mark recursive stage equations as global \verb|simp| lemmas.
\end{formalnote}

\section{Signs and finite coefficients}

\begin{implementation}[Sign type]
Use an inductive type, not \(\operatorname{Fin} 3\):
\begin{verbatim}
inductive Sign3 where
  | neg
  | zero
  | pos
  deriving DecidableEq, Repr

open Sign3

def sgn : Sign3 -> Int
  | neg  => -1
  | zero => 0
  | pos  => 1

def opp : Sign3 -> Sign3
  | neg  => pos
  | zero => zero
  | pos  => neg

def defect : Sign3 -> Nat
  | neg  => 2
  | zero => 1
  | pos  => 0
\end{verbatim}
The identity \((\defect(s):\Z)=1-\sgn(s)\) is the key bridge from signs to Nat-valued defect totals.
\end{implementation}

\begin{lemma}[Sign package]\label{lem:v9-sign-package}
For every \(s:\Sgn\), prove the following by cases:
\[
        \sgn(\opp(s))=-\sgn(s),
        \qquad \opp(\opp(s))=s,
\]
\[
        \nabs{\sgn(s)}\le1,
        \qquad (\defect(s):\Z)=1-\sgn(s),
\]
\[
        \defect(s)\le2,
        \qquad \sgn(s)=0\Longleftrightarrow s=\operatorname{zero}.
\]
\end{lemma}

\begin{definition}[Coefficient record]\label{def:v9-coeffon}
For \(A:\Finset\,\N\), define a coefficient record
\[
\CoeffOn(A)=\{\epsilon:\N\to\Sgn\;:\; \forall n, n\notin A\to \epsilon(n)=\operatorname{zero}\}.
\]
For \(c:\CoeffOn(A)\), write \(c(n)\) for its sign value.
\end{definition}

\begin{definition}[Grade, value, and signed difference]\label{def:v9-grade-value-diff}
Define
\[
        \SumN(A)=\sum_{a\in A}a\in\N,
\]
\[
        \Grade(A,c)=\sum_{a\in A}\sgn(c(a))\in\Z,
        \qquad
        \Val(A,c)=\sum_{a\in A}\sgn(c(a))(a:\Z)\in\Z.
\]
Then define
\[
        \Diff(A,t,x)\Longleftrightarrow
        \exists c:\CoeffOn(A),\ \Grade(A,c)=t\ \wedge\ \Val(A,c)=x.
\]
Suggested Lean shape:
\begin{verbatim}
structure CoeffOn (A : Finset Nat) where
  sign : Nat -> Sign3
  zero_of_not_mem : forall n, n notin A -> sign n = Sign3.zero

def grade (A : Finset Nat) (c : CoeffOn A) : Int :=
  A.sum (fun a => sgn (c.sign a))

def value (A : Finset Nat) (c : CoeffOn A) : Int :=
  A.sum (fun a => sgn (c.sign a) * (a : Int))

def Diff (A : Finset Nat) (t x : Int) : Prop :=
  exists c : CoeffOn A, grade A c = t /\ value A c = x
\end{verbatim}
\end{definition}

\begin{formalnote}[Why the support-zero field helps]
The support-zero field is not mathematically necessary, because all sums are over \(A\).  It is useful in Lean because restriction, union decomposition, and extensionality can be stated without repeatedly saying that coefficients outside the set are irrelevant.
\end{formalnote}

\begin{lemma}[Coefficient extensionality]\label{lem:v9-coeff-ext}
If \(c,d:\CoeffOn(A)\) satisfy
\[
        \forall a\in A,\ c(a)=d(a),
\]
then
\[
        \Grade(A,c)=\Grade(A,d),
        \qquad
        \Val(A,c)=\Val(A,d).
\]
\end{lemma}

\begin{lemma}[Negating coefficients]\label{lem:v9-coeff-neg}
For every \(c:\CoeffOn(A)\), define \(-c:\CoeffOn(A)\) by \((-c)(a)=\opp(c(a))\).  Then
\[
        \Grade(A,-c)=-\Grade(A,c),
        \qquad
        \Val(A,-c)=-\Val(A,c).
\]
Consequently
\[
        \Diff(A,t,x)\Rightarrow \Diff(A,-t,-x).
\]
\end{lemma}

\begin{lemma}[Finite bounds]\label{lem:v9-diff-bounds}
If \(\Diff(A,t,x)\), then
\[
        \nabs t\le \card A,
        \qquad
        \nabs x\le \SumN(A).
\]
The proof uses \(\nabs{\sum z_i}\le\sum\nabs{z_i}\), \(\nabs{\sgn(s)}\le1\), and \(\nabs{\sgn(s)(a:\Z)}\le a\).
\end{lemma}

\begin{definition}[Admissibility]\label{def:v9-admissibility}
Define the formal finite predicate by
\[
        \AdmF(A)\Longleftrightarrow
        \forall t:\Z,\ t\ne0\to\neg\Diff(A,t,0).
\]
Define the paper finite predicate by
\[
\begin{aligned}
\PaperAdmF(A)\Longleftrightarrow
\forall U,V:\Finset\,\N,
& U\subseteq A\to V\subseteq A\to\\
& \sum_{u\in U}u=\sum_{v\in V}v\to \card U=\card V.
\end{aligned}
\]
For an infinite set \(S:\Set\,\N\), define \(\PaperAdmS(S)\) similarly, using membership in \(S\).
\end{definition}

\begin{lemma}[Subset-pair bridge]\label{lem:v9-subset-bridge}
For every finite \(A:\Finset\,\N\),
\[
        \AdmF(A)\Longleftrightarrow \PaperAdmF(A).
\]
The two directions are proved by filtering a coefficient into positive and negative sign classes, and by building a coefficient from \(U\setminus V\) and \(V\setminus U\).
\end{lemma}

\section{Local block as range sums}

\begin{definition}[Local block data]\label{def:v9-local-block}
For \(m:\N\), define
\[
        \celt(m,i)=1+i(m+1),
        \qquad
        \Cblock(m)=\{\celt(m,i):0\le i<m\},
\]
\[
        \Qfun(m)=m^2+m+1.
\]
In Lean, define \(\Cblock(m)\) as the image of \(\operatorname{Finset.range} m\) under \(i\mapsto\celt(m,i)\).
\end{definition}

\begin{lemma}[Block arithmetic]\label{lem:v9-block-arithmetic}
For \(0<m\),
\[
        \card{\Cblock(m)}=m,
        \qquad
        \max \Cblock(m)=m^2,
        \qquad
        2\SumN(\Cblock(m))=m^3+m,
\]
\[
        \celt(m,i)\equiv1\pmod {m+1},
        \qquad
        \Qfun(m)=m(m+1)+1.
\]
The image-cardinality proof should use injectivity of \(i\mapsto1+i(m+1)\) on \(\operatorname{range}m\).
\end{lemma}

\begin{definition}[Range-local grade, index, and value]\label{def:v9-local-sums}
For \(m:\N\) and \(\epsilon:\N\to\Sgn\), define
\[
        \LGrade(m,\epsilon)=\sum_{i\in\range(m)}\sgn(\epsilon(i))\in\Z,
\]
\[
        \LIndex(m,\epsilon)=\sum_{i\in\range(m)}(i:\Z)\sgn(\epsilon(i))\in\Z,
\]
\[
        \LVal(m,\epsilon)=\sum_{i\in\range(m)}\sgn(\epsilon(i))(\celt(m,i):\Z)\in\Z.
\]
Then
\[
        \LVal(m,\epsilon)=\LGrade(m,\epsilon)+(m+1)\LIndex(m,\epsilon).
\]
Suggested Lean shape:
\begin{verbatim}
def localGrade (m : Nat) (eps : Nat -> Sign3) : Int :=
  (Finset.range m).sum (fun i => sgn (eps i))

def localIndex (m : Nat) (eps : Nat -> Sign3) : Int :=
  (Finset.range m).sum (fun i => (i : Int) * sgn (eps i))

def localValue (m : Nat) (eps : Nat -> Sign3) : Int :=
  (Finset.range m).sum
    (fun i => sgn (eps i) * ((1 + i * (m + 1) : Nat) : Int))
\end{verbatim}
\end{definition}

\begin{definition}[Local signed difference]\label{def:v9-cdiff}
Define
\[
        \CDiff(m,t,y)\Longleftrightarrow
        \exists \epsilon:\N\to\Sgn,
        \ \LGrade(m,\epsilon)=t\ \wedge\ \LVal(m,\epsilon)=y.
\]
This is the only local-difference predicate used in the local carry proof.
\end{definition}

\begin{lemma}[Local bounds]\label{lem:v9-cdiff-bounds}
If \(\CDiff(m,t,y)\), then
\[
        \nabs t\le m.
\]
If \(\epsilon\) is a witness, then
\[
        y=t+(m+1)\LIndex(m,\epsilon).
\]
\end{lemma}

\begin{lemma}[Block-to-local bridge]\label{lem:v9-block-to-local}
For \(0<m\),
\[
        \Diff(\Cblock(m),t,y)\Longleftrightarrow \CDiff(m,t,y).
\]
The forward direction chooses the sign attached to the unique index \(i<m\) with \(a=1+i(m+1)\).  The reverse direction pushes the range coefficient forward through the image map.
\end{lemma}

\section{Local carry: one canonical proof path}

\begin{definition}[Integer smallness predicates]\label{def:v9-local-small}
For \(m:\N\) and \(u,w,h:\Z\), define
\[
        \USmall(m,u)\Longleftrightarrow 4\nabs u<m,
        \qquad
        \WSmall(m,w)\Longleftrightarrow 2\nabs w<m,
\]
\[
        \HSmall(m,h)\Longleftrightarrow 4\nabs h<3m.
\]
These replace \(|u|<m/4\), \(|w|<m/2\), and \(|h|<3m/4\).  No rational or real division appears in the local proof.
\end{definition}

\begin{lemma}[Small subtraction]\label{lem:v9-small-sub}
If \(\USmall(m,u)\) and \(\WSmall(m,w)\), then
\[
        \HSmall(m,w-u).
\]
This is the triangle inequality for \(\operatorname{Int.natAbs}\), multiplied by \(4\).
\end{lemma}

\begin{lemma}[Local quotient identity]\label{lem:v9-local-quotient}
Let \(0<m\).  Suppose \(\epsilon:\N\to\Sgn\),
\[
        t=\LGrade(m,\epsilon),
        \qquad
        y=\LVal(m,\epsilon),
        \qquad
        y=\Qfun(m)w-u.
\]
Then there is \(\ell:\Z\) such that
\[
        t=(w-u)+\ell(m+1).
\]
Proof normal form: rewrite \(y=t+(m+1)I\) and \(\Qfun(m)w-u=(w-u)+(m+1)mw\), then take \(\ell=mw-I\).
\end{lemma}

\begin{lemma}[Three residue cases]\label{lem:v9-three-residue-cases}
Let \(0<m\).  Suppose
\[
        \nabs t\le m,
        \qquad
        \HSmall(m,h),
        \qquad
        t=h+\ell(m+1).
\]
Then
\[
        \ell=-1\quad\vee\quad \ell=0\quad\vee\quad \ell=1.
\]
This lemma should be a standalone integer arithmetic theorem.  It should not mention blocks, coefficients, or finite sums.
\end{lemma}

\begin{corollary}[Three local alternatives]\label{cor:v9-three-local-alternatives}
Let \(0<m\).  Suppose
\[
        \CDiff(m,t,y),
        \qquad y=\Qfun(m)w-u,
        \qquad \USmall(m,u),
        \qquad \WSmall(m,w).
\]
Then
\[
        t=w-u
        \quad\vee\quad
        t=w-u+(m+1)
        \quad\vee\quad
        t=w-u-(m+1).
\]
\end{corollary}

\subsection{Defect totals}

\begin{definition}[Defect totals]\label{def:v9-defect-totals}
For \(m:\N\) and \(\epsilon:\N\to\Sgn\), define
\[
        \DTotal(m,\epsilon)=\sum_{i\in\range(m)}\defect(\epsilon(i))\in\N,
\]
\[
        \WDTotal(m,\epsilon)=\sum_{i\in\range(m)}i\defect(\epsilon(i))\in\N.
\]
\end{definition}

\begin{lemma}[Defect-grade bridge]\label{lem:v9-defect-grade}
For every \(m\) and \(\epsilon\),
\[
        (\DTotal(m,\epsilon):\Z)=(m:\Z)-\LGrade(m,\epsilon).
\]
This is a direct sum of \((\defect(s):\Z)=1-\sgn(s)\).
\end{lemma}

\begin{lemma}[Doubled index-defect identity]\label{lem:v9-doubled-index-defect}
For every \(m\) and \(\epsilon\),
\[
        2\LIndex(m,\epsilon)+2(\WDTotal(m,\epsilon):\Z)
        =(m:\Z)((m:\Z)-1).
\]
This avoids division by \(2\) and keeps the identity compatible with \verb|ring|.
\end{lemma}

\begin{lemma}[Weighted defect bound]\label{lem:v9-weighted-defect-bound}
If \(0<m\), then
\[
        \WDTotal(m,\epsilon)\le (m-1)\DTotal(m,\epsilon).
\]
Use \(i\le m-1\) for \(i\in\operatorname{range}m\), multiply by the nonnegative defect, and sum.
\end{lemma}

\subsection{Positive shift contradiction}

\begin{lemma}[Positive shift extracts the defect total]\label{lem:v9-pos-shift-defect}
Let \(0<m\), let \(\epsilon:\N\to\Sgn\), and put
\[
        t=\LGrade(m,\epsilon),
        \qquad
        R=\DTotal(m,\epsilon).
\]
If
\[
        t=w-u+(m+1),
\]
then
\[
        u=w+(R:\Z)+1.
\]
This is the cast-safe replacement for the paper's auxiliary integer \(a=u-w\).
\end{lemma}

\begin{lemma}[Positive shift index identity]\label{lem:v9-pos-shift-index}
Let \(0<m\), let \(\epsilon:\N\to\Sgn\), and put
\[
        t=\LGrade(m,\epsilon),
        \qquad
        y=\LVal(m,\epsilon),
        \qquad
        I=\LIndex(m,\epsilon).
\]
If
\[
        y=\Qfun(m)w-u,
        \qquad
        t=w-u+(m+1),
\]
then
\[
        I=(m:\Z)w-1.
\]
\end{lemma}

\begin{lemma}[Positive shift lower bound]\label{lem:v9-pos-shift-lower}
Let \(0<m\), let \(\epsilon:\N\to\Sgn\), and put
\[
        I=\LIndex(m,\epsilon),
        \qquad R=\DTotal(m,\epsilon).
\]
Then
\[
        4I\ge
        2(m:\Z)((m:\Z)-1)-4((m:\Z)-1)(R:\Z).
\]
Use Lemmas~\ref{lem:v9-doubled-index-defect} and \ref{lem:v9-weighted-defect-bound}.
\end{lemma}

\begin{lemma}[Positive shift arithmetic core]\label{lem:v9-pos-shift-arith-core}
Let \(0<m\).  Let \(R:\N\) and \(I,u,w:\Z\).  Suppose
\[
        I=(m:\Z)w-1,
        \qquad
        u=w+(R:\Z)+1,
        \qquad
        \USmall(m,u),
\]
\[
        4I\ge
        2(m:\Z)((m:\Z)-1)-4((m:\Z)-1)(R:\Z).
\]
Then \(\operatorname{False}\).
After substituting the first two equalities, the contradiction is pure arithmetic.  The normalized positive remainder is
\[
        (m:\Z)^2+2(m:\Z)+4(R:\Z)+4>0.
\]
This should be the only dense \verb|omega|/\verb|nlinarith| lemma in the local carry file.
\end{lemma}

\begin{lemma}[Positive shifted case impossible]\label{lem:v9-pos-shift-impossible}
Let \(0<m\).  Suppose
\[
        \CDiff(m,t,y),
        \qquad y=\Qfun(m)w-u,
        \qquad \USmall(m,u),
        \qquad t=w-u+(m+1).
\]
Then \(\operatorname{False}\).
Proof route: choose the \(\epsilon\) witnessing \(\CDiff\), set \(R=\DTotal(m,\epsilon)\), apply Lemmas~\ref{lem:v9-pos-shift-defect}, \ref{lem:v9-pos-shift-index}, \ref{lem:v9-pos-shift-lower}, and \ref{lem:v9-pos-shift-arith-core}.
\end{lemma}

\subsection{Negative shift by sign negation}

\begin{lemma}[Negating local coefficients]\label{lem:v9-local-neg}
For \(\epsilon:\N\to\Sgn\), define \(\epsilon^{-}(i)=\opp(\epsilon(i))\).  Then
\[
        \LGrade(m,\epsilon^{-})=-\LGrade(m,\epsilon),
        \qquad
        \LIndex(m,\epsilon^{-})=-\LIndex(m,\epsilon),
        \qquad
        \LVal(m,\epsilon^{-})=-\LVal(m,\epsilon).
\]
Consequently
\[
        \CDiff(m,t,y)\Rightarrow\CDiff(m,-t,-y).
\]
\end{lemma}

\begin{lemma}[Negative shifted case impossible]\label{lem:v9-neg-shift-impossible}
Let \(0<m\).  Suppose
\[
        \CDiff(m,t,y),
        \qquad y=\Qfun(m)w-u,
        \qquad \USmall(m,u),
        \qquad t=w-u-(m+1).
\]
Then \(\operatorname{False}\).
Apply Lemma~\ref{lem:v9-pos-shift-impossible} to \((-t,-y,-w,-u)\).
\end{lemma}

\begin{theorem}[Local carry, range form]\label{thm:v9-local-carry-range}
Let \(0<m\).  Suppose
\[
        \CDiff(m,t,y),
        \qquad y=\Qfun(m)w-u,
        \qquad \USmall(m,u),
        \qquad \WSmall(m,w).
\]
Then
\[
        t=w-u.
\]
Use Corollary~\ref{cor:v9-three-local-alternatives}; the two shifted cases are Lemmas~\ref{lem:v9-pos-shift-impossible} and \ref{lem:v9-neg-shift-impossible}.
\end{theorem}

\begin{corollary}[Local carry for the finite block]\label{cor:v9-local-carry-block}
Let \(0<m\).  Suppose
\[
        \Diff(\Cblock(m),t,y),
        \qquad y=\Qfun(m)w-u,
        \qquad \USmall(m,u),
        \qquad \WSmall(m,w).
\]
Then \(t=w-u\).  This is Theorem~\ref{thm:v9-local-carry-range} plus Lemma~\ref{lem:v9-block-to-local}.
\end{corollary}

\section{Integer-only carry closure}

\begin{definition}[Smallness predicates]\label{def:v9-smallness}
For \(\sigma,M,k:\N\), define
\[
        \OldSmall(\sigma,M,k)\Longleftrightarrow 4\sigma<Mk,
        \qquad
        \NewSmall(\sigma,M,k)\Longleftrightarrow 2\sigma<Mk.
\]
These are the only smallness predicates in the finite construction.
\end{definition}

\begin{lemma}[Quotient smallness]\label{lem:v9-quotient-smallness}
Let \(0<M\).  If
\[
        D=(M:\Z)q,
        \qquad
        \nabs D\le \sigma,
        \qquad
        \OldSmall(\sigma,M,k),
\]
then \(\USmall(k,q)\).  If instead \(\NewSmall(\sigma,M,k)\), then \(\WSmall(k,q)\).
The proof uses
\[
        \nabs{(M:\Z)q}=M\nabs q
\]
and cancellation of the positive natural factor \(M\).
\end{lemma}

\begin{definition}[Cached carry state]\label{def:v9-carry-state}
A carry state is a record
\[
\CarryState(A,M,\sigma)
\]
with fields
\[
\begin{array}{ll}
\text{Mpos} & 0<M,\\
\text{sigmaEq} & \sigma=\SumN(A),\\
\text{posMem} & \forall a\in A,\ 0<a,\\
\text{ltM} & \forall a\in A,\ a<M,\\
\text{carry} & \forall T,D,q:\Z,\ \Diff(A,T,D)\to D=(M:\Z)q\to T=q.
\end{array}
\]
The cached field \(\sigma\) prevents repeated unfolding of \(\SumN(A)\) in closure and stage induction.
\end{definition}

\begin{lemma}[Carry state implies admissible]\label{lem:v9-state-admissible}
If \(\CarryState(A,M,\sigma)\), then \(\AdmF(A)\) and \(\PaperAdmF(A)\).
\end{lemma}

\begin{definition}[Legal step]\label{def:v9-legal-step}
For a state \(s=\CarryState(A,M,\sigma)\), define
\[
        \LegalStep(s,k)\Longleftrightarrow 0<k\ \wedge\ \OldSmall(\sigma,M,k).
\]
The constructed next block is
\[
        \block(M,k)=M\Cblock(k),
\]
with next modulus
\[
        M'=M\Qfun(k)
\]
and cached next sum
\[
        \sigma'=\sigma+M\SumN(\Cblock(k)).
\]
\end{definition}

\begin{lemma}[Block separation]\label{lem:v9-block-separation}
If \(s=\CarryState(A,M,\sigma)\) and \(0<k\), then
\[
        A\cap\block(M,k)=\varnothing,
        \qquad
        \forall b\in\block(M,k),\ M\le b<M\Qfun(k).
\]
The first statement follows from \(a<M\) for old elements and \(M\le b\) for block elements.
\end{lemma}

\begin{lemma}[Dilation bridge]\label{lem:v9-dilation-bridge}
Let \(0<M\) and \(0<k\).  If
\[
        \Diff(\block(M,k),t,x),
\]
then there is \(y:\Z\) such that
\[
        \CDiff(k,t,y),
        \qquad
        x=(M:\Z)y.
\]
This is the finite-image bridge for the dilation \(c\mapsto M c\).  It should be proved once and reused only in closure.
\end{lemma}

\begin{lemma}[Disjoint-union decomposition]\label{lem:v9-union-decompose}
If \(A\cap B=\varnothing\) and
\[
        \Diff(A\cup B,T,D),
\]
then there exist \(T_A,T_B,D_A,D_B:\Z\) such that
\[
        \Diff(A,T_A,D_A),
        \qquad
        \Diff(B,T_B,D_B),
\]
\[
        T=T_A+T_B,
        \qquad
        D=D_A+D_B.
\]
The proof restricts the coefficient record to \(A\) and \(B\).  The support-zero field in \(\CoeffOn\) is designed for this lemma.
\end{lemma}

\begin{lemma}[Post-step smallness]\label{lem:v9-post-step-smallness}
Let \(0<M\), \(0<k\), and
\[
        Q=\Qfun(k),
        \qquad
        \sigma'=\sigma+M\SumN(\Cblock(k)).
\]
If
\[
        2\SumN(\Cblock(k))=k^3+k,
        \qquad
        \OldSmall(\sigma,M,k),
\]
then
\[
        \NewSmall(\sigma',MQ,k).
\]
Arithmetic normal form after rewriting the doubled block sum:
\[
        2\sigma+M(k^3+k)<M(k^3+k^2+k).
\]
From \(4\sigma<Mk\) and \(0<k\), derive \(2\sigma<Mk^2\); the rest is \(\N\)-arithmetic.
\end{lemma}

\begin{theorem}[One-step closure, cached-state form]\label{thm:v9-closure-step}
Let \(s=\CarryState(A,M,\sigma)\), and suppose \(\LegalStep(s,k)\).  Put
\[
        A'=A\cup\block(M,k),
        \qquad
        M'=M\Qfun(k),
        \qquad
        \sigma'=\sigma+M\SumN(\Cblock(k)).
\]
Then
\[
        \CarryState(A',M',\sigma')
        \qquad\text{and}\qquad
        \NewSmall(\sigma',M',k).
\]
\end{theorem}

\begin{proof}[Formal proof route]
Only the carry field is nontrivial.  Suppose
\[
        \Diff(A',T,D),
        \qquad D=(M':\Z)w.
\]
Use Lemma~\ref{lem:v9-union-decompose} and Lemma~\ref{lem:v9-dilation-bridge} to write
\[
        T=T_A+T_B,
        \qquad
        D=D_A+(M:\Z)y,
        \qquad
        \Diff(A,T_A,D_A),
        \qquad
        \CDiff(k,T_B,y).
\]
Since \(D=(M:\Z)\Qfun(k)w\), define
\[
        u=\Qfun(k)w-y.
\]
Then
\[
        D_A=(M:\Z)u.
\]
The old carry field gives \(T_A=u\).  Lemma~\ref{lem:v9-quotient-smallness}, using old smallness and \(\nabs{D_A}\le\sigma\), gives \(\USmall(k,u)\).  Lemma~\ref{lem:v9-post-step-smallness} gives \(\NewSmall(\sigma',M',k)\), and Lemma~\ref{lem:v9-quotient-smallness}, using \(\nabs D\le\sigma'\), gives \(\WSmall(k,w)\).  The local carry theorem gives
\[
        T_B=w-u.
\]
Therefore
\[
        T=T_A+T_B=u+(w-u)=w.
\]
\end{proof}

\begin{lemma}[New smallness becomes next old smallness]\label{lem:v9-new-to-old}
If
\[
        \NewSmall(\sigma,M,k_{old}),
        \qquad
        k_{old}\le N,
        \qquad
        2N<k_{new},
\]
then
\[
        \OldSmall(\sigma,M,k_{new}).
\]
Indeed,
\[
        4\sigma<2Mk_{old}\le2MN<Mk_{new}.
\]
\end{lemma}

\section{Stage construction}

\begin{definition}[Abstract schedule]\label{def:v9-schedule}
A schedule consists of functions
\[
        k,N,M:\N\to\N
\]
with recursive equations
\[
        N_0=0,
        \qquad M_0=1,
        \qquad N_{j+1}=N_j+k_j,
        \qquad M_{j+1}=M_j\Qfun(k_j),
\]
and stage hypotheses
\[
        0<k_j,
        \qquad
        2N_{j+1}<k_{j+1}\quad\text{for all }j.
\]
The concrete ceiling schedule is introduced only in Section~\ref{sec:v9-concrete-ceiling}.
\end{definition}

\begin{definition}[Finite stages]\label{def:v9-finite-stages}
Given a schedule, define
\[
        F_0=\varnothing,
        \qquad
        B_j=\block(M_j,k_j),
        \qquad
        F_{j+1}=F_j\cup B_j.
\]
Also define
\[
        \sigma_j=\SumN(F_j).
\]
For implementation, it is often better to define \(\sigma_j\) recursively by the cached-sum formula and prove \(\sigma_j=\SumN(F_j)\) as a field of the carry state.
\end{definition}

\begin{theorem}[Stage invariant]\label{thm:v9-stage-invariant}
For every \(j:\N\),
\[
        \CarryState(F_j,M_j,\sigma_j).
\]
For every \(j\), after adding block \(j\),
\[
        \NewSmall(\sigma_{j+1},M_{j+1},k_j).
\]
For every \(j\),
\[
        \card{F_j}=N_j.
\]
\end{theorem}

\begin{proof}[Induction route]
The base state is \(F_0=\varnothing\), \(M_0=1\), \(\sigma_0=0\).  The base legal condition for \(k_0\) is trivial because \(4\sigma_0=0<k_0M_0\).  The step from \(j\) to \(j+1\) is Theorem~\ref{thm:v9-closure-step}.  To make the next legal condition, use Lemma~\ref{lem:v9-new-to-old}, the fact \(k_j\le N_{j+1}\), and the schedule hypothesis \(2N_{j+1}<k_{j+1}\).  Cardinality follows from block separation and \(\card{\Cblock(k_j)}=k_j\).
\end{proof}

\begin{corollary}[Admissibility of all finite stages]\label{cor:v9-finite-stage-adm}
For every \(j\), \(F_j\) is admissible in both the formal and paper senses.
\end{corollary}

\begin{definition}[Infinite set]\label{def:v9-infinite-set}
Define
\[
        A_\infty=\{x:\N:\exists j,\ x\in B_j\}.
\]
Equivalently, \(x\in A_\infty\) iff \(x\in F_j\) for some \(j\).
\end{definition}

\begin{theorem}[Infinite admissibility]\label{thm:v9-infinite-adm}
The set \(A_\infty\) is paper-admissible.
\end{theorem}

\begin{proof}
If finite subsets \(U,V\subset A_\infty\) have equal sums, choose a stage \(J\) containing every element of \(U\cup V\).  Apply Corollary~\ref{cor:v9-finite-stage-adm} to \(F_J\).
\end{proof}

\section{Exact enumeration without sorting the infinite set}

\begin{definition}[Block element]\label{def:v9-block-element}
For \(j,i:\N\), define
\[
        \elem(j,i)=M_j\bigl(1+i(k_j+1)\bigr).
\]
For \(i<k_j\), this is the \(i\)-th element of block \(B_j\), using zero-based indexing within the block.
\end{definition}

\begin{lemma}[Order inside and between blocks]\label{lem:v9-block-order}
For \(i+1<k_j\),
\[
        \elem(j,i)<\elem(j,i+1).
\]
For every \(j\),
\[
        \elem(j,k_j-1)<\elem(j+1,0)
\]
when \(0<k_j\).  More precisely,
\[
        \elem(j+1,0)-\elem(j,k_j-1)=M_j(k_j+1).
\]
For \(i<j\), every element of \(B_i\) is smaller than every element of \(B_j\).
\end{lemma}

\begin{definition}[Owner and offset]\label{def:v9-owner-offset}
For \(n:\N\), define \(\owner(n)\) to be the unique \(j\) such that
\[
        N_j\le n<N_{j+1}.
\]
Use \(\operatorname{Nat.find}\) on existence.  Define
\[
        \offset(n)=n-N_{\owner(n)}.
\]
Then
\[
        \offset(n)<k_{\owner(n)}.
\]
\end{definition}

\begin{definition}[Zero-based sequence]\label{def:v9-seq}
Define
\[
        \seq(n)=\elem(\owner(n),\offset(n)).
\]
This is the zero-based increasing enumeration.  The paper sequence is \(a_{n+1}=\seq(n)\).
\end{definition}

\begin{theorem}[Exact block formula]\label{thm:v9-exact-block-formula}
If \(i<k_j\), then
\[
        \seq(N_j+i)=\elem(j,i).
\]
This theorem should be the main way to rewrite the sequence; do not unfold \(\owner\) directly in later files.
\end{theorem}

\begin{theorem}[Uniform gap formula]\label{thm:v9-uniform-gap}
If \(i<k_j\), then
\[
        \seq(N_j+i+1)-\seq(N_j+i)=M_j(k_j+1).
\]
The theorem is intentionally uniform: the same formula covers the internal gaps \(i+1<k_j\) and the inter-block gap \(i+1=k_j\).
\end{theorem}

\begin{lemma}[Owner tends to infinity]\label{lem:v9-owner-atTop}
\[
        \owner(n)\to\atTop\qquad(n\to\atTop).
\]
It suffices to prove that for every fixed \(J\), all indices \(n\ge N_{J+1}\) have \(\owner(n)>J\).  This uses monotonicity of \(N_j\) and positivity of every \(k_j\).
\end{lemma}

\section{Abstract asymptotic API}\label{sec:v9-asymptotic-api}

\begin{definition}[Parameters]\label{def:v9-params}
Set
\[
        \thetaC=1+\sqrt2,
        \qquad
        \tauC=2+\sqrt2,
        \qquad
        \fexp=3+2\sqrt2.
\]
Then
\[
        \fexp=\thetaC+\tauC,
        \qquad
        \tauC=\thetaC\sqrt2,
        \qquad
        \tauC+\thetaC(2-\tauC)=0,
        \qquad
        2-\tauC=-\sqrt2.
\]
These algebraic facts should be separate lemmas near the start of the real-analysis file.
\end{definition}

\begin{definition}[Scaled modulus]\label{def:v9-scaled-modulus}
Use the Version 10 total scale
\[
        \Scale(j)=\max\{1,N_j\},
        \qquad
        \Pfun(j)=\frac{(M_j:\R)}{\rpow((\Scale(j):\R),\tauC)}.
\]
This avoids the denominator \(N_0=0\) without shifting indices.  All block-ratio lemmas should use \(\Pfun(\owner(n))\), not \(P_{\owner(n)-1}\).
\end{definition}

\begin{lemma}[Geometric decay of \(\Pfun\)]\label{lem:v9-P-geometric}
For the concrete ceiling schedule, there is a constant
\[
        \rho=2^{1-\sqrt2}
\]
with \(0\le\rho<1\) such that, for every \(j\),
\[
        \Pfun(j+1)\le \rho \Pfun(j).
\]
The proof normal form is
\[
\begin{aligned}
\Pfun(j+1)
&=\frac{M_{j+1}}{\Scale(j+1)^{\tauC}}
\le \frac{2M_jk_j^2}{k_j^{\tauC}}
=2M_jk_j^{2-\tauC} \\
&\le 2M_j(2\Scale(j)^{\thetaC})^{-\sqrt2}
=2^{1-\sqrt2}\frac{M_j}{\Scale(j)^{\thetaC\sqrt2}}
=\rho \Pfun(j).
\end{aligned}
\]
Here \(\Scale(j+1)\ge k_j\), and the lower bound for \(k_j\) is stated in terms of \(\Scale(j)\).  No \(j\ge1\) split is needed.
\end{lemma}

\begin{lemma}[Geometric decay tends to zero]\label{lem:v9-geom-to-zero}
Let \(P:\N\to\R\), \(0\le\rho\), and \(\rho<1\).  If \(0\le P_j\) and
\[
        P_{j+1}\le\rho P_j
\]
eventually or for all \(j\), then
\[
        P_j\to0.
\]
Use this as a general real-analysis lemma.  The admissible-set proof should not reprove convergence of a geometrically dominated sequence.
\end{lemma}

\begin{lemma}[Ceiling upper bound for block length]\label{lem:v9-k-upper}
There is a constant \(C_0>0\) such that, for every \(j\),
\[
        k_j+1\le C_0\Scale(j)^{\thetaC}.
\]
For \(j=0\), \(\Scale(0)=1\) and \(k_0=2\), so increasing \(C_0\) once handles the base case.  For \(j>0\), use the usual ceiling upper bound and \(N_j+1\le2N_j\).
\end{lemma}

\begin{lemma}[Block gap ratio bound]\label{lem:v9-block-gap-ratio}
For every block \(j\) and every \(i<k_j\), put \(n=N_j+i\).  Then
\[
        \frac{(\seq(n+1)-\seq(n):\R)}{\rpow((n+1:\R),\fexp)}
        \le C_0 \Pfun(j).
\]
Use Theorem~\ref{thm:v9-uniform-gap}, \(n+1\ge \Scale(j)\), \(\fexp=\thetaC+\tauC\), and Lemma~\ref{lem:v9-k-upper}.  The base block is included because \(\Scale(0)=1\).
\end{lemma}

\begin{lemma}[Owner-bound transfer]\label{lem:v9-owner-bound-transfer}
Let \(R:\N\to\R\) and \(B:\N\to\R\) be nonnegative eventually.  Suppose
\[
        \owner(n)\to\atTop,
        \qquad
        B_j\to0,
\]
and eventually
\[
        R(n)\le C B_{\owner(n)}
\]
for some constant \(C\).  Then
\[
        R(n)\to0.
\]
This isolates the filter bookkeeping in the final gap proof.
\end{lemma}

\begin{theorem}[Abstract gap little-oh theorem]\label{thm:v9-abstract-gap-little-o}
Assume the stage, enumeration, geometric decay, and block-length upper-bound lemmas above.  Then
\[
        \frac{(\seq(n+1)-\seq(n):\R)}{(n+1:\R)^{\fexp}}\to0
        \qquad(n\to\infty).
\]
Equivalently, for the one-based paper enumeration \(a_{n+1}=\seq(n)\),
\[
        a_{n+1}-a_n=o(n^{3+2\sqrt2}).
\]
\end{theorem}

\section{Concrete ceiling schedule}\label{sec:v9-concrete-ceiling}

\begin{definition}[Concrete schedule]\label{def:v9-concrete-schedule}
Use the recursive equations
\[
        N_0=0,
        \qquad M_0=1,
\]
\[
        k_j=\left\lceil 2(1+N_j)^{\thetaC}\right\rceil,
        \qquad
        N_{j+1}=N_j+k_j,
        \qquad
        M_{j+1}=M_j\Qfun(k_j).
\]
In Lean, define \(k_j\) using a natural ceiling.  The real expression should be converted once into a natural number, and all later finite-stage proofs should use only the derived inequalities listed below.
\end{definition}

\begin{lemma}[Concrete schedule finite inequalities]\label{lem:v9-concrete-finite-ineq}
For the concrete schedule, prove the following \(\N\)-level facts:
\[
        0<k_j,
        \qquad
        k_0=2,
        \qquad
        2N_{j+1}<k_{j+1},
        \qquad
        k_j\le N_{j+1}.
\]
Only the proof of \(2N_{j+1}<k_{j+1}\) should look at \(\operatorname{Real.rpow}\).  After this lemma, the stage induction is purely finite.
\end{lemma}

\begin{lemma}[Concrete schedule real inequalities]\label{lem:v9-concrete-real-ineq}
For every \(j\), prove
\[
        k_j\ge 2\Scale(j)^{\thetaC},
        \qquad
        Q_{k_j}\le 2k_j^2,
        \qquad
        \Scale(j+1)\ge k_j,
\]
and the upper estimate in Lemma~\ref{lem:v9-k-upper}.  These are the only concrete growth facts needed by Lemma~\ref{lem:v9-P-geometric} and the final gap theorem.  The constant in the upper estimate may be enlarged once to cover \(j=0\).
\end{lemma}

\begin{theorem}[Final theorem, Lean target]\label{thm:v9-final}
For the concrete ceiling schedule, the infinite set
\[
        A_\infty=\bigcup_{j\ge0}B_j
\]
is paper-admissible, infinite, has first block \(B_0=\{1,4\}\), and its one-based increasing enumeration \(a_1<a_2<\cdots\) satisfies
\[
        a_{n+1}-a_n=o\bigl(n^{3+2\sqrt2}\bigr).
\]
Consequently
\[
        a_{n+1}-a_n\le n^{3+2\sqrt2}
\]
for all sufficiently large \(n\).
\end{theorem}

\section{Minimal Lean skeleton}

\begin{contract}[Core declarations]
The following skeleton is intentionally close to Lean syntax but not meant to be pasted without adapting names.
\begin{verbatim}
inductive Sign3 where | neg | zero | pos

def sgn : Sign3 -> Int := ...
def opp : Sign3 -> Sign3 := ...
def defect : Sign3 -> Nat := ...

structure CoeffOn (A : Finset Nat) where
  sign : Nat -> Sign3
  zero_of_not_mem : forall n, n notin A -> sign n = Sign3.zero

def sumN (A : Finset Nat) : Nat := A.sum id

def grade (A : Finset Nat) (c : CoeffOn A) : Int := ...
def value (A : Finset Nat) (c : CoeffOn A) : Int := ...
def Diff (A : Finset Nat) (t x : Int) : Prop := ...

def cell (m i : Nat) : Nat := 1 + i * (m + 1)
def C (m : Nat) : Finset Nat := (Finset.range m).image ...
def Q (m : Nat) : Nat := m*m + m + 1

def localGrade (m : Nat) (eps : Nat -> Sign3) : Int := ...
def localIndex (m : Nat) (eps : Nat -> Sign3) : Int := ...
def localValue (m : Nat) (eps : Nat -> Sign3) : Int := ...
def CDiff (m : Nat) (t y : Int) : Prop := ...

def defectTotal (m : Nat) (eps : Nat -> Sign3) : Nat := ...
def weightedDefectTotal (m : Nat) (eps : Nat -> Sign3) : Nat := ...

def USmall (m : Nat) (u : Int) : Prop := 4 * u.natAbs < m
def WSmall (m : Nat) (w : Int) : Prop := 2 * w.natAbs < m

theorem local_carry_zsmall {m : Nat} (hm : 2 <= m)
    (eps : Nat -> Sign3) {u w : Int}
    (hy : localValue m eps = (Q m : Int) * w - u)
    (hu : ZLtQuarter m u)
    (hw : ZLtHalf m w) :
    localGrade m eps = w - u := ...

theorem local_carry_range {m : Nat} (hm : 0 < m)
    {t y u w : Int}
    (hcd : CDiff m t y)
    (hy : y = (Q m : Int) * w - u)
    (hu : USmall m u)
    (hw : WSmall m w) :
    t = w - u := ...

structure CarryState (A : Finset Nat) (M sigma : Nat) where
  Mpos : 0 < M
  sigmaEq : sigma = sumN A
  posMem : forall a, a in A -> 0 < a
  ltM : forall a, a in A -> a < M
  carry : forall T D q : Int, Diff A T D -> D = (M : Int) * q -> T = q

def OldSmall (sigma M k : Nat) : Prop := 4 * sigma < M * k
def NewSmall (sigma M k : Nat) : Prop := 2 * sigma < M * k

theorem closure_step ... :
  CarryState A' M' sigma' /\ NewSmall sigma' M' k := ...

def elem (j i : Nat) : Nat := M j * (1 + i * (k j + 1))
def owner (n : Nat) : Nat := Nat.find ...
def offset (n : Nat) : Nat := n - N (owner n)
def seq (n : Nat) : Nat := elem (owner n) (offset n)

theorem uniform_gap {j i : Nat} (hi : i < k j) :
  seq (N j + i + 1) - seq (N j + i) = M j * (k j + 1) := ...

def AdmSeq (a : Nat -> Nat) : Prop :=
  forall U V : Finset Nat, U.sum a = V.sum a -> U.card = V.card

def PrefixRange (a : Nat -> Nat) (F : Nat -> Finset Nat) (N : Nat -> Nat) : Prop :=
  forall j, F j = (Finset.range (N j)).image a
\end{verbatim}
\end{contract}

\section{Version 10 proof-obligation ledger}

\begin{longtable}{>{\raggedright\arraybackslash}p{0.32\linewidth} >{\raggedright\arraybackslash}p{0.46\linewidth} >{\raggedright\arraybackslash}p{0.14\linewidth}}
Obligation & Main tools & Difficulty \\ \hline
Sign package & \verb|cases|, \verb|simp| & 1/10 \\
Coefficient extensionality and negation & \verb|Finset.sum_congr|, sign package & 2/10 \\
Subset-pair bridge & filters on finite sets, extensionality & 3/10 \\
Diff bounds & finite-sum triangle inequality, \verb|Int.natAbs| & 3/10 \\
Block arithmetic for \(C_k\) & image injectivity, sum over range & 3/10 \\
Range local value identity & \verb|Finset.sum_congr|, \verb|ring| & 2/10 \\
Three residue cases & isolated integer arithmetic & 3/10 \\
Defect-grade and doubled index identities & range sums, sign package, \verb|ring| & 3/10 \\
Weighted defect bound & monotonicity over \(\operatorname{range}m\) & 3/10 \\
Positive shift arithmetic core & \verb|omega|/\verb|nlinarith| after normalization & 4/10 \\
Two-sided local carry theorem & z-small wrappers plus shift contradictions & 3.75/10 \\
Dilation bridge & image of finite set under multiplication & 4/10 \\
Disjoint-union decomposition & coefficient restriction & 4/10 \\
Closure step & quotient extraction plus local carry & 4.5/10 \\
Stage invariant & induction, cached records & 3.5/10 \\
Owner and exact sequence formula & \verb|Nat.find| and monotone \(N_j\) & 4/10 \\
Uniform gap formula & internal/boundary split, same result & 3.5/10 \\
Geometric decay of \(\Pfun\) & max-one scale, \(\operatorname{Real.rpow}\), positivity, constants & 4.25/10 \\
Owner-bound transfer & filters, squeeze theorem & 4/10 \\
Concrete final theorem & assembly only & 3/10
\end{longtable}

\section{Recommended implementation order}

\begin{checklist}[Linearized order]
\begin{enumerate}[label=(\arabic*), leftmargin=*]
\item Implement \(\Sgn\), \(\sgn\), \(\opp\), \(\defect\), and prove Lemma~\ref{lem:v9-sign-package}.
\item Implement \(\CoeffOn\), \(\Grade\), \(\Val\), \(\Diff\), coefficient extensionality, coefficient negation, and bounds.
\item Prove the subset-pair bridge and define admissibility predicates.
\item Implement \(\celt\), \(\Cblock\), \(\Qfun\), local range sums, and \(\CDiff\).
\item Prove the block-to-local bridge.
\item Prove local-carry arithmetic in this exact order: NatAbs firewall, two-sided subtraction bound, shift identity, three-shift arithmetic core, defect-grade bridge, doubled index-defect identity, weighted defect bound, positive shift extraction, positive shift index identity, positive shift lower bound, positive shift arithmetic core, positive shift impossible, local negation, negative shift impossible, final local carry wrapper.
\item Implement integer smallness predicates and quotient smallness.
\item Implement cached carry states, block separation, dilation bridge, disjoint-union decomposition, post-step smallness, and closure.
\item Implement abstract schedule and stage induction.
\item Define block owner, offset, and zero-based sequence; prove exact block formula and uniform gap formula.
\item Prove the abstract asymptotic API with \(\Scale(j)=\max\{1,N_j\}\): geometric decay to zero, owner-bound transfer, and abstract gap little-oh.
\item Instantiate the concrete ceiling schedule and assemble the final theorem.
\end{enumerate}
\end{checklist}

\begin{formalnote}[What to avoid]
Avoid the following patterns, because they usually make the Lean proof longer.
\begin{enumerate}[label=(\roman*), leftmargin=*]
\item Do not prove local carry directly from subset witnesses.
\item Do not use real absolute values or rational division in the finite files.
\item Do not unfold the infinite set and sort it to define the enumeration.
\item Do not keep several local-carry proof paths alive in the same Lean file.
\item Do not let the concrete ceiling formula appear in the closure or enumeration files.
\end{enumerate}
\end{formalnote}

\end{document}
